\metatitle{Pólya frequency sequences and total positivity}

\metadescription{Pólya frequency sequences, totally non-negative matrices, Toeplitz matrices, Edrei--Thoma theory, and real-rootedness applications.}


\section[polyaFrequency]{Pólya frequency sequences}

Pólya frequency sequences give a coefficient-level way to recognize
real-rooted polynomials with non-negative coefficients.  The basic idea is to
put the coefficients into a Toeplitz matrix and ask for total
non-negativity.  This links real-rootedness with planar networks,
\hyperref[interlacingPolynomials]{interlacing polynomials}, and
\hyperref[hadamardProduct]{Hadamard products}.

\subsection[totallyNonnegative]{Totally non-negative matrices}

A matrix \(M=(m_{ij})\) is \defin{totally non-negative} (TNN) if every minor
is non-negative.  It is \defin{totally positive} (TP) if every minor is
strictly positive.

TNN matrices arise naturally from path-counting in directed graphs via the
\hyperref[lgvLemma]{Lindström--Gessel--Viennot lemma} (LGV).  Planar networks
with compatible source and sink orderings yield TNN matrices automatically.
As a standard example, Pascal's triangle \(M_{n,j}=\binom{n}{j}\) is TNN by
the LGV lemma applied to a grid with West and North steps.
See \cite{ChenFuRuan2025x} for more on this viewpoint.

\begin{example*}[TNN for a matrix via LGV]
The matrix \(M\) defined by
\[
  M_{n,j} = \binom{d(n-j\ell)}{j}
\]
admits a lattice path interpretation.  The entry \(M_{n,j}\) counts
West--North paths from \(A_n=(dn,0)\) to
\[
  B_j=((d\ell+1)j,j).
\]
Since sources lie on the \(x\)-axis with spacing \(d\), and sinks lie on the
line \(y=x/(d\ell+1)\), any non-identity matching forces a path crossing.
By the \hyperref[lgvLemma]{LGV lemma}, the matrix \((M_{n,j})\) is TNN.
\end{example*}

\subsection[toeplitzTnnMatrix]{Toeplitz matrices and PF sequences}

\begin{definition}[Pólya frequency sequence]\label{def:polyaFrequencySequence}
A sequence \(a=(a_0,a_1,\dotsc)\) of non-negative real numbers is a
\defin{Pólya frequency sequence} (PF sequence) if the semi-infinite Toeplitz
matrix
\[
  T(a)_{i,j}=a_{i-j},\qquad a_k=0 \text{ for } k<0,
\]
is totally non-negative.
\end{definition}

Thus a PF sequence is a sequence whose finite Toeplitz minors are all
non-negative.  In particular, PF sequences are log-concave and have no
internal zeros.

\begin{theorem}[Aissen--Schoenberg--Whitney, 1952]
\label{aissenSchoenbergWhitney}
A polynomial
\[
  p(t)=\sum_k a_k t^k,\qquad a_k\geq 0,
\]
has only real non-positive roots if and only if its coefficient sequence
\((a_0,a_1,\dotsc)\) is a Pólya frequency sequence.
See \cite{Karlin1968} for a full treatment.
\end{theorem}

The theorem is often the fastest way to move between real-rootedness and
matrix total positivity.  For example, if a construction takes a TNN Toeplitz
matrix to another TNN Toeplitz matrix, it automatically preserves
real-rootedness of the corresponding polynomial.

\subsection[edreiThomaTheorem]{Edrei--Thoma theorem}

The infinite-sequence version has a precise analytic form.

\begin{theorem}[Edrei--Thoma theorem, see \cite{Karlin1968}]
A non-negative sequence \((a_n)_{n\geq 0}\) is a Pólya frequency sequence if
and only if its ordinary generating function has the form
\[
  \sum_{n\geq 0} a_n x^n
  =
  Cx^m e^{\gamma x}
  \prod_i \frac{1+\alpha_i x}{1-\beta_i x},
\]
where \(C,\gamma,\alpha_i,\beta_i\geq 0\), \(m\geq 0\), and
\(\sum_i(\alpha_i+\beta_i)<\infty\).
\end{theorem}

The polynomial case is the finite part of this classification: a polynomial
with non-negative coefficients has coefficient sequence PF exactly when its
zeros are real and non-positive.

\subsection[veroneseSectionsPF]{Veronese sections}

\begin{theorem}[Veronese sections preserve real-rootedness]
\label{veroneseSectionsRealRooted}
Let \(r\geq 1\) and \(0\leq i<r\).  Define the \defin{Veronese section}
\[
  \mathcal{S}^r_i\left(\sum_{m\geq 0} a_m t^m\right)
  =
  \sum_{m\geq 0} a_{i+rm} t^m.
\]
If \(f(t)\) is real-rooted with non-negative coefficients, then
\(\mathcal{S}^r_i(f)\) is also real-rooted with non-negative coefficients.
\end{theorem}

\begin{proof}
The roots of \(f\) are necessarily non-positive.  Hence, by
\hyperref[aissenSchoenbergWhitney]{Aissen--Schoenberg--Whitney}, the
coefficient sequence \((a_0,a_1,\dotsc)\) is PF, so its Toeplitz matrix
\[
  T_{p,q}=a_{p-q}
\]
is TNN.  The Toeplitz matrix for the subsequence
\((a_i,a_{i+r},a_{i+2r},\dotsc)\) is the submatrix
\[
  T_{i+rp,rq}=a_{i+r(p-q)}.
\]
It is therefore TNN as well.  Applying Aissen--Schoenberg--Whitney again
shows that \(\mathcal{S}^r_i(f)\) is real-rooted.
\end{proof}

\name{Christos A. Athanasiadis} and \name{David G. Wagner} prove a stronger
\hyperref[fullyInterlacingLace]{fully interlacing} version of this statement
\cite[Thm.~5.3 and Rem.~5.4]{AthanasiadisWagner2024x}.  In particular, if
\(f\) has non-negative coefficients and only real roots, then
\[
  \bigl(
    \mathcal{S}^r_0(f),
    \mathcal{S}^r_1(f),
    \dotsc,
    \mathcal{S}^r_{r-1}(f)
  \bigr)
\]
is fully interlacing.  Thus the Veronese construction produces an interlacing
family of polynomials, not only separate real-rooted polynomials.

\begin{example*}[Rows of Pascal's triangle]
The polynomial
\[
  \sum_{m=0}^{n}\binom{kn}{km}t^m
\]
is the Veronese section \(\mathcal{S}^k_0((1+t)^{kn})\).
Since \((1+t)^{kn}\) is real-rooted, the theorem above immediately gives
that this polynomial is real-rooted.
\end{example*}

\subsection[tnnPathMatrix]{Path matrices and total positivity}

\begin{theorem}[Characterization of TNN matrices, \cite[Thm.~3.1]{Brenti1995}]
\label{brentiTnnPathMatrix}
A matrix is TNN if and only if it can be obtained from a fully compatible
planar directed graph with non-negative edge weights.
\end{theorem}

\begin{theorem}[Characterization of TNN lower unitriangular matrices, \cite{BrandenLeite2024x}]
\label{tnnLowerUnitriangularMatrix}
Let \(\Gamma_N\) be the directed graph on
\(\{(i,j)\in\mathbb{N}^2 : j \leq i \leq N\}\) with horizontal edges
\((i,j)\to(i,j+1)\) and vertical edges \((i+1,j)\to(i,j)\).  Assign
non-negative weights
\(\lambda=(\lambda_{i,j})_{0\leq j\leq i<N}\) to the vertical edges
(and weight \(1\) on horizontal edges), and let
\(r_{n,k}(\Gamma_N,\lambda)\) be the total weight of paths from \((n,0)\)
to \((k,k)\).

A lower unitriangular matrix \(R=(r_{n,k})_{n,k=0}^N\) is totally
non-negative if and only if there exists such a weight array \(\lambda\) with
\(R=R(\Gamma_N,\lambda)\).  The forward direction follows from the
\hyperref[lgvLemma]{LGV lemma}; the converse is Whitney's reduction theorem,
see \cite{FallatJohnson2011}.
\end{theorem}

\emph{Checking TNN.}
Total non-negativity can be checked efficiently via \emph{Neville
elimination}, a variant of Gaussian elimination using adjacent-row
operations.  A matrix is TNN if and only if Neville elimination completes
with all multipliers non-negative and no entry becoming negative.  This gives
an \(O(n^2m)\) algorithm for an \(n\times m\) matrix, compared to the
exponential cost of checking all minors directly.

TNN matrices are closed under products and under Hadamard products, see
\cite{Wagner1992}.  See \cite{BrandenLeite2024x} for applications related to
chains in posets, and \cite{FallatJohnson2011} for a comprehensive treatment.

\subsection[polyaFrequencyExamples]{Examples}

\begin{example*}[Ray through Pascal's triangle]
The following is proved in \cite{Yu2009}.  Let \(n\geq k\geq 0\) and
\(b>a>0\), with \(k<b\).  Then the polynomial
\[
  \sum_{j\geq 0} \binom{n+ja}{k+jb} t^j
\]
is real-rooted.  The proof uses an interpretation via lattice paths and
solves \cite[Problem~12]{ForgcsTran2016}.
\end{example*}

\begin{example*}[Delannoy numbers]
Let \(\Delta(a,b)\) be the \defin{Delannoy number} counting lattice paths
from \((0,0)\) to \((a,b)\) with east, north, and north-east steps.
Equivalently,
\[
  \Delta(a,b)=\sum_i \binom{a}{i}\binom{b}{i}2^i .
\]
The square array \((\Delta(a,b))_{a,b\geq 0}\) is \oeis{A008288} when read
by antidiagonals, and it forms a TNN matrix \cite[Cor.~5.15]{Brenti1995}.
The finite antidiagonal row polynomials
\[
  P_n(t)=\sum_{k=0}^n \Delta(n-k,k)t^k
\]
begin
\[
  1,\qquad 1+t,\qquad 1+3t+t^2,\qquad 1+5t+5t^2+t^3,\dotsc
\]
and satisfy
\[
  P_0(t)=1,\qquad P_1(t)=1+t,\qquad
  P_n(t)=(1+t)P_{n-1}(t)+tP_{n-2}(t) \quad (n\geq 2).
\]
These are the \hyperref[delannoyRowPolynomials]{Delannoy row polynomials}
that occur as \(h^*\)-polynomials of regular snake posets.
\end{example*}

\begin{theorem}[Brenti~\cite{Brenti1995}]\label{thm:brenti}
Let \(D\) be a locally finite digraph on \(\setN\times\setN\) with
non-negative edge weights.  Assume that \(D\) is planar and weakly
\(y\)-invariant, meaning that the edges from \((m,k)\) are independent
of \(m\).  Let
\[
  M_{n,j}\coloneqq P_D\bigl((0,0),(n,j)\bigr)
\]
be the weighted path count from \((0,0)\) to \((n,j)\).  Then
\(M=(M_{n,j})\) is totally positive and every row is a Pólya frequency
sequence.  In particular, every row polynomial \(\sum_j M_{n,j}t^j\)
has only real non-positive zeros.
\end{theorem}

\begin{theorem}[Stanley, \cite{Stanley1998}]
\label{stanleySchurPositiveRealRooted}
A polynomial \(Q(x)=1+a_1x+a_2x^2+\dotsb+a_nx^n\) has only negative real
roots if and only if the product
\[
  \prod_{i\geq 1} Q(x_i)
\]
is Schur-positive.
\end{theorem}

\name{Ziyao Fu}, \name{Yulin Peng}, and \name{Yuchong Zhang} connect this
kind of total positivity with upper homogeneous posets, or \defin{upho
posets} \cite{FuPengZhang2024x}.  They prove that a formal power series
\(f(x)\in 1+x\setZ_{\geq 0}[[x]]\) is the rank-generating function of a
Schur-positive upho poset if and only if \(f(x)\) is totally positive.  Here
Schur positivity refers to the Ehrenborg quasisymmetric function of the poset.
Their proof passes through a correspondence between colored upho posets and
atomic left-cancellative monoids.
